\color{unidarkgreen}{\subsection{Максимальная токовая защита}}
\color{black}

\subsubsection{Максимальная токовая защиты стороны ВН}

\begin{enumerate}[label=\arabic{section}.\arabic{subsection}.\arabic{subsubsection}.\arabic*,  align=left, labelsep=4pt, leftmargin=0pt, itemindent=57pt]

\item
Максимальная токовая защита (МТЗ) стороны ВН предназначена для резервирования действия
основных быстродействующих защит. Защита выполнена двухступенчатой. Измерительные цепи
защиты включены на токи фаз А, В, С со стороны ВН трансформатора.    

\item
Каждая из ступеней имеет регулируемые значения по току и времени срабатывания. Значения
срабатывания по току определяются параметрами \textbf{«Iср МТЗ-1 ВН»} и \textbf{«Iср МТЗ-2 ВН»} для первой и
второй ступеней соответственно. Задержки на срабатывание защиты задаются параметрами \textbf{«Тср МТЗ
-1 ВН»} и \textbf{«Тср МТЗ-2 ВН»}.

\item
Каждая из ступеней защиты может выполнена с пуском по напряжению. Ввод контроля пуска по
напряжению в логику работы ступеней осуществляется программными накладками \textbf{«Режим КПОН МТЗ-
1 ВН»} и \textbf{«Режим КПОН МТЗ-2 ВН»} для первой и второй ступеней соответственно. 

\item
Каждая из ступеней защиты может выполнена с блокировкой по второй гармонике (отношению
уровня второй гармоники к первой). Данная блокировка реализована на базе модуля «ОВ БНТ ВН». Ввод
контроля броска намагничивающего тока программными накладками  \textbf{«Режим БНТ МТЗ-1 ВН»}, \textbf{«Режим
БНТ МТЗ-2 ВН»}.

\item
Для обеспечения несрабатывания МТЗ ВН при броске тока намагничивания используется
перекрестная блокировка по 2 гармонике в фазных токах ВН. Вводится программной накладкой \textbf{«БНТ
МТЗ ВН»}. Блокировка основана на контроле отношения действующего значения второй гармонике к
первой токов ВН. Данная блокировка вводиться в работу для каждой ступени независимо с помощью
уставок \textbf{«Режим БНТ МТЗ-1 ВН»} и \textbf{«Режим БНТ МТЗ-2 ВН»}. Имеется возможность регулировать порог
срабатывания блокировки, с помощью уставки \textbf{«I2г/I1г МТЗ ВН»}. Блокировка МТЗ ВН по 2 гармонике
выводится из действия при величине соответствующего фазного тока ниже 0,04Iном. Также блокировка
МТЗ ВН по 2 гармонике выводится из действия при достижении величины действующего значения тока
выше уставки \textbf{«Iср БНТ МТЗ ВН»}.

\begin{figure}[H]
\centering
\includegraphics[width=0.8\textwidth,height=0.8\textheight,keepaspectratio]{\fbpath/021(022,023).0311(0321, 0331) МТЗ/img/MTZ_VN.pdf}
\captionof{figure}{Функциональный блок МТЗ ВН}\label{pic:MTZ_VN}
\end{figure}

\end{enumerate}


\subsubsection{Максимальная токовая защиты стороны НН1(НН)}

\begin{enumerate}[label=\arabic{section}.\arabic{subsection}.\arabic{subsubsection}.\arabic*,  align=left, labelsep=4pt, leftmargin=0pt, itemindent=57pt]

\item
Максимальная токовая защита (МТЗ) стороны НН1(НН) предназначена для резервирования
действия основных быстродействующих защит. Защита выполнена двухступенчатой. Измерительные
цепи защиты включены на токи фаз А, В, С со стороны НН1(НН) трансформатора. 

\item
Каждая из ступеней имеет регулируемые значения по току и времени срабатывания. Значения срабатывания по
току определяются параметрами \textbf{«Iср МТЗ-1 НН1(НН)»} и \textbf{«Iср МТЗ-2 НН1(НН)»} для
первой и второй ступеней соответственно. Задержки на срабатывание защиты задаются параметрами
\textbf{«Тср МТЗ-1 НН1(НН)»} и \textbf{«Тср МТЗ-2 НН1(НН)»}.

\item
Каждая из ступеней защиты может выполнена с пуском по напряжению. Ввод контроля пуска по
напряжению в логику работы ступеней осуществляется программными накладками \textbf{«Режим КПОН МТЗ-
1 НН1(НН)»} и \textbf{«Режим КПОН МТЗ-2 НН1(НН)»} для первой и второй ступеней соответственно.

\item
Каждая из ступеней защиты может выполнена с блокировкой по второй гармонике (отношению
уровня второй гармоники к первой). Данная блокировка реализована на базе модуля «ОВ БНТ НН1(НН)».
Ввод контроля броска намагничивающего тока программными накладками \textbf{«Режим БНТ МТЗ-1
НН1(НН)»}, \textbf{«Режим БНТ МТЗ-2 НН1(НН)»}. 

\item
Для обеспечения несрабатывания МТЗ НН1(НН) при броске тока намагничивания используется
перекрестная блокировка по 2 гармонике в фазных токах НН1(НН). Вводится программной накладкой
\textbf{«БНТ МТЗ НН1(НН)»}.  Блокировка основана на контроле отношения действующего значения второй
гармонике к первой токов НН1(НН). Данная блокировка вводиться в работу для каждой ступени
независимо с помощью уставок \textbf{«Режим БНТ МТЗ-1 НН1(НН)»} и \textbf{«Режим БНТ МТЗ-2 НН1(НН)»}. Имеется
возможность регулировать порог срабатывания блокировки, с помощью уставки \textbf{«I2г/I1г МТЗ НН1(НН)»}.
Блокировка МТЗ НН1(НН) по 2 гармонике выводится из действия при величине соответствующего
фазного тока ниже 0,04Iном. Также блокировка МТЗ НН1(НН) по 2 гармонике выводится из действия при
достижении величины действующего значения тока выше уставки \textbf{«Iср БНТ МТЗ НН1(НН)»}.

\begin{figure}[H]
\centering
\includegraphics[width=0.8\textwidth,height=0.8\textheight,keepaspectratio]{\fbpath/021(022,023).0311(0321, 0331) МТЗ/img/MTZ_NN1.pdf}
\captionof{figure}{Функциональный блок МТЗ НН1(НН)}\label{pic:MTZ_NN1}
\end{figure}

\end{enumerate}

\subsubsection{Максимальная токовая защиты стороны НН2(CН)}

\begin{enumerate}[label=\arabic{section}.\arabic{subsection}.\arabic{subsubsection}.\arabic*,  align=left, labelsep=4pt, leftmargin=0pt, itemindent=57pt]

\item
Направленная максимальная токовая защита на стороне НН2(СН) с комбинированным пуском по
напряжению предназначена для резервирования основных защит трансформатора, резервирования
отключения КЗ на шинах НН2(СН) и на элементах присоединенных к шинам НН2(СН).

\begin{figure}[H]
\centering
\includegraphics[width=0.8\textwidth,height=0.8\textheight,keepaspectratio]{\fbpath/021(022,023).0311(0321, 0331) МТЗ/img/MTZ_NN2.pdf}
\captionof{figure}{Функциональный блок МТЗ НН2(СН)}\label{pic:MTZ_NN2}
\end{figure}

\item
Каждая из ступеней имеет регулируемые значения по току и времени срабатывания. Значения
срабатывания по току определяются параметрами \textbf{«Iср МТЗ-1 НН2(СН)»} и \textbf{«Iср МТЗ-2 НН2(СН)»} для
первой и второй ступеней соответственно. Задержки на срабатывание защиты задаются параметрами
\textbf{«Тср МТЗ-1 НН2(СН)»} и \textbf{«Тср МТЗ-2 НН2(СН)»}.

\item
Каждая из ступеней защиты может выполнена с пуском по напряжению. Ввод контроля пуска по
напряжению в логику работы ступеней осуществляется программными накладками \textbf{«Режим КПОН МТЗ
-1 НН2(СН)»} и \textbf{«Режим КПОН МТЗ-2 НН2(СН)»} для первой и второй ступеней соответственно.

\item
Каждая из ступеней защиты может выполнена с блокировкой по второй гармонике (отношению
уровня второй гармоники к первой). Данная блокировка реализована на базе модуля «ОВ БНТ НН2(СН)».
Ввод контроля броска намагничивающего тока программными накладками \textbf{«Режим БНТ МТЗ-1
НН2(СН)»}, \textbf{«Режим БНТ МТЗ-2 НН2(СН)»}.

\item
Для обеспечения несрабатывания МТЗ НН2(СН) при броске тока намагничивания используется
перекрестная блокировка по 2 гармонике в фазных токах НН2(СН). Вводится программной накладкой
\textbf{«БНТ МТЗ НН2(СН)»}. Блокировка основана на контроле отношения действующего значения второй
гармонике к первой токов НН2(СН). Данная блокировка вводиться в работу для каждой ступени
независимо с помощью уставок \textbf{«Режим БНТ МТЗ-1 НН2(СН)»} и \textbf{«Режим БНТ МТЗ-2 М»}. Имеется
возможность регулировать порог срабатывания блокировки, с помощью уставки \textbf{«I2г/I1г МТЗ НН2(СН)»}.
Блокировка МТЗ НН2(СН) по 2 гармонике выводится из действия при величине соответствующего
фазного тока ниже 0,04Iном. Также блокировка МТЗ НН2(СН) по 2 гармонике выводится из действия при
достижении величины действующего значения тока выше уставки \textbf{«Iср БНТ МТЗ НН2(СН)»}.

\medskip

\begin{figure}[H]
\centering
\begin{subfigure}{0.45\textwidth}
    \centering
    \includegraphics[width=\textwidth, trim=0 10mm 150mm 0, clip]{\fbpath/021(022,023).0311(0321, 0331) МТЗ/img/ONM.pdf}
    \caption{а) ОНМ включен на Ia и Ubc}
    \label{pict:ONM_a}
\end{subfigure}
%\hfill
\begin{subfigure}{0.45\textwidth}
    \centering
    \includegraphics[width=\textwidth, trim=150mm 10mm 0 0, clip]{\fbpath/021(022,023).0311(0321, 0331) МТЗ/img/ONM.pdf}
    \caption{б) ОНМ включен на Ic и Uab}
    \label{pict:ONM_b}
\end{subfigure}

\caption{Функциональная схема ОНМ МТЗ СН Т}
\label{pict:ONM}
\end{figure}

\medskip

\item
Орган направления мощности (рисунок \ref{pict:ONM}) выполнен по 90-градусной схеме сочетания токов и
напряжений Ia и Ubc, Ic и Uab. Направление мощности определяется по величине фазового угла между
током и напряжением отдельно для каждой пары сигналов. Разрешение работы МТЗ НН2(СН) будет
происходить по паре сигналов ОНМ с максимальным током. Углы максимальной чувствительности ОНМ
в сторону трансформатора и в сторону сети задаются уставками \textbf{«Фмч ОНМ – В Т»} и \textbf{«Фмч ОНМ – В
сеть»}. Угол максимальной чувствительности откладывается от линейного напряжения, за
положительное направление принято вращение против часовой стрелки. Для ОНМ направленного в
трансформатор, рекомендуется выбирать угол максимальной чувствительности в диапазоне 15º…60º.
Для ОНМ направленного в сеть, рекомендуется выбирать угол максимальной чувствительности в
диапазоне 195º…250º.  Величина тока работы ОНМ 0,05*Iном, величина напряжения работы ОНМ
0,01*Uном.

\item
Режим направленности для каждой ступени задается с помощью уставок \textbf{«Направ МТЗ-1
НН2(СН)»} и \textbf{«Направ МТЗ-2 НН2(СН)»}:
\begin{itemize}
\item \textit{«Ненапр»} - МТЗ СН без учета направления;
\item \textit{«В Т»} - МТЗ СН направлена в сторону трансформатора;
\item \textit{«В сеть»} - МТЗ СН  направлена в сторону сети.
\end{itemize}

\item
Режим работы направленных ступеней МТЗ НН2(СН) при неисправности ТН задается с помощью
уставки \textbf{«КЦН ОНМ МТЗ»}:
\begin{itemize}
\item \textit{«Выведено»} - Режим работы ступени МТЗ НН2(СН) не зависит от БНН;
\item \textit{«РНМ»} - Ступень МТЗ НН2(СН) переводится в ненаправленный режим работы;
\item \textit{«МТЗ»} - Вывод направленных ступеней МТЗ НН2(СН).
\end{itemize}

\item
Формирование сигнала срабатывания КПОН низкой и средней стороны трансформатора
производится с контролем положения выключателей соответствующих сторон.

\small
\begin{longtable}{|>{\centering\arraybackslash}m{5.3cm}|>{\centering\arraybackslash}m{4.2cm}|>{\centering\arraybackslash}m{1cm}|>{\centering\arraybackslash}m{3.3cm}|>{\centering\arraybackslash}m{1.8cm}|}
% Первая страница: полный заголовок
\caption{Перечень уставочных параметров МТЗ\hfill\vspace{-0.5\baselineskip}}\label{desc:tbl1}\\
\hline
\rowcolor{gray!30}
Наименование уставки & Диапазон & Шаг & Значение по умолчанию & Единица измерения \\
\hline
\endfirsthead

% Промежуточные страницы: «Продолжение таблицы»
\caption*{\hspace{3pt}\emph{Продолжение таблицы \ref{desc:tbl1}}\hfill\ }\vspace{-0.5\baselineskip} \\
\hline
\rowcolor{gray!30}
Наименование уставки & Диапазон & Шаг & Значение по умолчанию & Единица измерения \\
\hline
\endhead

%===m>ДТЗ|-
%===m
\end{longtable}
\normalsize	

\end{enumerate}